% ==========================================================
%  main.tex  --  Manuscrito anonimizado (revisión doble ciega)
%  Clase: sagej.cls (SAGE)   Compilar: pdflatex -> bibtex -> pdflatex x2
%  Coloca sagej.cls y SageH.bst (de la plantilla SAGE) en esta carpeta.
% ==========================================================
\documentclass[Afour,sageh,times]{sagej}
% Opciones: sageh = Harvard (autor-año); usa sagev si la revista requiere Vancouver.
% Consulta siempre las normas de envío de la revista; prevalecen sobre la plantilla.

\usepackage[T1]{fontenc}
\usepackage[utf8]{inputenc}
\usepackage{xcolor}
\usepackage{xurl}
\usepackage[colorlinks,bookmarksopen,bookmarksnumbered,citecolor=black,urlcolor=black,linkcolor=black]{hyperref}

% Marcador provisional (se muestra en rojo hasta que se reemplace)
\newcommand{\relleno}[1]{\textcolor{red}{[#1]}}
\renewcommand{\refname}{Referencias}

\def\volumeyear{2026}

\begin{document}

\runninghead{Gobernanza y prevención de la corrupción en el Perú}

% --- Mantener nombres y afiliaciones únicamente en portada.tex. ---
\author{}  % Vacío intencionalmente para la revisión doble ciega
\title{Gobernanza y prevención de la corrupción en el Perú: desafíos institucionales, políticas de integridad y mecanismos de control en la gestión pública}

\begin{abstract}
Corruption prevention in Peru is closely linked to governance quality, state institutional capacity, and the coordination of integrity policies and control mechanisms. This article presents a narrative documentary review examining these relationships in the Peruvian context through five selected documents: the OECD's \textit{Integrity Review of Peru} (2017), \textit{Justice Review of Peru} (2024), and \textit{Towards a National Integrity and Transparency System in Peru} (2024); Eloy Munive Pariona's proposal for a public integrity system (2022); and Peru's National Policy on Integrity and the Fight against Corruption (2017). The review proceeds from the premise that public integrity depends not only on rules or specialised institutions, but also on their ability to coordinate, implement policies, manage risks, ensure accountability, and activate control and sanction mechanisms. It proposes an analytical framework with three dimensions: institutional coordination and capacity; integrity and prevention policies; and control, accountability, and sanction mechanisms. This framework connects preventive action with the organisational capacities required for integrity policies to have practical effects in Peruvian public management.
\end{abstract}

\keywords{public governance, corruption prevention, public integrity, institutional capacity, control and accountability, Peru}

\maketitle

%% ----------------------------------------------------------
\section{Introducción}
\subsection{Contexto y justificación}
La integridad pública es central para la calidad de la gobernanza porque condiciona la forma en que las instituciones ejercen sus funciones, administran recursos y responden a la ciudadanía. En el Perú, \citet{oecd2017integrity} sostiene que la integridad del sector público es importante para el desarrollo socioeconómico sostenido y que debe fortalecerse mediante una combinación de cultura de integridad, rendición de cuentas, control y aplicación efectiva de las normas. Su enfoque abarca los niveles nacional y subnacional e incluye ética pública, conflictos de intereses, protección de denunciantes, control interno, gestión de riesgos, sistemas disciplinarios y justicia penal.

Esta perspectiva es especialmente relevante porque la prevención de la corrupción no puede reducirse a identificar y sancionar posteriormente las conductas indebidas. La Política Nacional de Integridad y Lucha contra la Corrupción del Perú se organiza en torno a tres ejes: capacidad preventiva del Estado; identificación y gestión de riesgos; y capacidad sancionadora del Estado. Así, la política conecta prevención, gestión institucional y sanción como componentes de una misma respuesta estatal (\citealp{pcm2017politica}).

La evolución institucional posterior refuerza este problema. La revisión de la OCDE de 2024 informa que el Perú ha creado organismos especializados, desarrollado un Modelo de Integridad e incorporado Oficinas de Integridad Institucional y funcionarios responsables del acceso a la información. Sin embargo, esta arquitectura no garantiza por sí sola resultados homogéneos. La revisión identifica dificultades de coordinación interinstitucional, ambigüedad en las responsabilidades y limitaciones de recursos y apoyo institucional para implementar políticas de integridad y transparencia (\citealp{oecd2024sistema}).

Por tanto, estudiar la corrupción desde una perspectiva de gobernanza requiere examinar no solo las normas existentes, sino también las capacidades estatales necesarias para aplicarlas y coordinarlas. Esta relación ayuda a explicar por qué la prevención de la corrupción depende de instituciones capaces de actuar conjuntamente, gestionar riesgos y sostener mecanismos de control y rendición de cuentas (\citealp{oecd2017integrity,oecd2024sistema}).

\subsection{Vacío en la literatura}
La literatura seleccionada ofrece perspectivas diferenciadas sobre integridad, prevención, control, justicia y coordinación institucional en el Perú. La revisión de la OCDE de 2017 analiza el sistema de integridad y sus principales componentes de control y cumplimiento; la Política Nacional de Integridad y Lucha contra la Corrupción establece objetivos y ejes de política pública; y una propuesta académica plantea un Sistema Nacional de Integridad orientado a fortalecer la capacidad preventiva del Estado (\citealp{oecd2017integrity,pcm2017politica,munive2022integridad}).

La revisión de justicia de la OCDE también muestra que los problemas de coordinación y la fragmentación institucional afectan el funcionamiento del sistema de justicia peruano. Identifica enfoques compartimentados, duplicidad de funciones y dificultades para coordinar responsabilidades entre instituciones. Sostiene que la implementación de reformas depende de una estructura institucional que permita coordinar, planificar y supervisar las acciones de manera sistémica (\citealp{oecd2024justice}).

El vacío que aborda este artículo no consiste en la ausencia de estudios sobre corrupción o integridad en el Perú. Más bien, reúne tres dimensiones estrechamente relacionadas que requieren un análisis conjunto: capacidades institucionales y mecanismos de coordinación; políticas de integridad y prevención; y mecanismos de control, rendición de cuentas y sanción. Esta perspectiva integrada supera una visión fragmentada de instrumentos y organismos, y considera la prevención de la corrupción como una función sistémica de la gobernanza pública peruana.

\subsection{Objetivo y pregunta de investigación}
El objetivo general es analizar, mediante una revisión documental narrativa, la relación entre gobernanza, integridad pública, capacidad institucional, políticas de prevención y mecanismos de control frente a la corrupción en el Perú.

De manera específica, el artículo examina los desafíos institucionales asociados con la coordinación y las capacidades estatales; identifica los principales componentes de las políticas peruanas de integridad y prevención; analiza el papel del control, la rendición de cuentas y la sanción; y propone un marco analítico que relaciona estas dimensiones desde una perspectiva de gobernanza pública.

La revisión se guía por la siguiente pregunta de investigación: \textbf{¿Cómo se articulan la coordinación y las capacidades institucionales, las políticas de integridad y prevención, y los mecanismos de control, rendición de cuentas y sanción en la prevención de la corrupción en la gestión pública peruana?}

\subsection{Contribución y estructura del artículo}
La contribución del artículo consiste en organizar las fuentes seleccionadas desde una perspectiva integradora de gobernanza e integridad pública. En lugar de considerar la prevención, el control y la sanción como ámbitos independientes, los analiza como componentes interdependientes de la capacidad institucional necesaria para prevenir y responder a los riesgos de corrupción. Este enfoque es consistente con la propuesta de un sistema organizado en niveles estratégico, táctico y operativo, y con la propuesta de fortalecer la rectoría del sistema y la coordinación horizontal y vertical (\citealp{munive2022integridad,oecd2024sistema}).

El artículo se organiza en tres grandes partes. Primero, desarrolla la revisión de la literatura y el marco analítico. Segundo, presenta la metodología de revisión documental. Finalmente, presenta los hallazgos y la discusión derivados del análisis de las fuentes.

%% ----------------------------------------------------------
\section{Revisión de la literatura y marco analítico}
\subsection{Corrupción, gobernanza e integridad pública}
Las fuentes revisadas abordan la corrupción desde una perspectiva institucional y la relacionan con la capacidad del Estado para salvaguardar la integridad en el ejercicio de las funciones públicas. La integridad pública supone orientar la actuación institucional y la conducta de los servidores públicos por principios y valores compatibles con la protección del interés público, y respaldarlas con mecanismos que reduzcan las oportunidades de comportamientos contrarios a la ética. \citet{oecd2017integrity} vincula explícitamente integridad, gobernanza, control y gestión de riesgos, mientras que \citet{munive2022integridad} considera la integridad un componente fundamental de la prevención de la corrupción.

La relación con la gobernanza surge de la necesidad de implementar estas políticas de manera coordinada dentro de un sistema institucional. \citet{munive2022integridad} propone una estructura con niveles estratégico, táctico y operativo que incorpora actores estatales y no estatales. Por su parte, la revisión de la OCDE de 2024 sostiene que la arquitectura institucional debe facilitar la articulación entre las distintas entidades y niveles de gobierno, en lugar de limitar la integridad y la transparencia a mecanismos formales de cumplimiento (\citealp{oecd2024sistema}).

\subsection{Desafíos institucionales en la gestión pública peruana}
La coordinación es un desafío central. La existencia de múltiples entidades, funciones y mecanismos exige responsabilidades claras, canales de cooperación y capacidades de dirección y seguimiento. La revisión de la OCDE de 2024 identifica problemas de coordinación entre las entidades responsables de integridad y transparencia, así como ambigüedad en las responsabilidades y limitaciones de recursos y apoyo administrativo (\citealp{oecd2024sistema}).

La capacidad institucional es un segundo aspecto. Las políticas de integridad deben adaptarse a distintas realidades organizacionales y territoriales. La revisión de la OCDE de 2024 señala que trasladar los estándares de integridad y transparencia a entidades de diferentes tamaños, funciones y capacidades constituye un desafío para implementar las políticas de manera coherente (\citealp{oecd2024sistema}).

Un tercer desafío corresponde a la relación entre control y gestión. La revisión de integridad de la OCDE sobre el Perú señala que el control interno y la gestión de riesgos pueden tratarse como rutinas administrativas aisladas e integrarse de manera insuficiente en los procesos de gestión. Por ello, la capacidad preventiva depende de que los responsables de la administración incorporen la gestión de riesgos y el control en las decisiones y procesos institucionales (\citealp{oecd2017integrity}).

\subsection{Políticas de integridad y anticorrupción}
La Política Nacional de Integridad y Lucha contra la Corrupción del Perú es el principal referente de política pública seleccionado para esta revisión. Se estructura en torno a tres ejes: capacidad preventiva; identificación y gestión de riesgos; y capacidad sancionadora. El primero comprende transparencia, gestión de la información, cultura de integridad y gestión de conflictos de intereses; el segundo abarca supervisión e identificación y gestión de riesgos; y el tercero busca fortalecer la capacidad estatal para investigar, perseguir y sancionar la corrupción (\citealp{pcm2017politica}).

\citet{munive2022integridad} amplía esta perspectiva al proponer un Sistema Nacional de Integridad capaz de articular actores públicos y privados en distintos niveles de intervención. Esta lógica sistémica conecta la formulación de políticas con su implementación institucional y busca fortalecer la capacidad preventiva de las entidades públicas.

\citet{oecd2024sistema} retoma esta problemática y propone una mayor coherencia institucional mediante un Sistema Nacional de Integridad y Transparencia, una rectoría fortalecida y mecanismos de coordinación interinstitucional. Este planteamiento muestra que una política de integridad requiere no solo instrumentos, sino también una arquitectura que defina cómo se relacionan los actores responsables de su aplicación.

\subsection{Mecanismos de prevención y control}
Los mecanismos de prevención incluyen transparencia, gestión de conflictos de intereses, identificación de riesgos, sistemas de denuncia, integridad organizacional y control interno. Estos instrumentos buscan actuar antes de que las prácticas corruptas se materialicen o consoliden. En particular, la Política Nacional vincula la identificación y gestión de riesgos con mecanismos permanentes de supervisión, mientras que la revisión de la OCDE de 2017 relaciona el control interno con la gestión institucional de riesgos (\citealp{pcm2017politica,oecd2017integrity}).

Los mecanismos de control también deben relacionarse con la rendición de cuentas y la asignación de responsabilidades. La Política Nacional incorpora explícitamente la capacidad sancionadora como uno de sus ejes, mientras que la revisión de la OCDE de 2017 considera los sistemas disciplinarios y el papel de la justicia penal. En el ámbito de la justicia, la revisión de la OCDE de 2024 subraya que la coordinación institucional, la claridad de responsabilidades y la capacidad de implementación son condiciones pertinentes para que los mecanismos de control funcionen eficazmente (\citealp{pcm2017politica,oecd2017integrity,oecd2024justice}).

\subsection{Marco analítico}
La revisión propone un marco analítico compuesto por tres dimensiones interrelacionadas.

\textbf{Primera dimensión: coordinación y capacidades institucionales.} Comprende la claridad de competencias, la coordinación horizontal y vertical, los recursos disponibles, la profesionalización, la capacidad de implementación y la existencia de una entidad rectora capaz de orientar y supervisar la política de integridad. Esta dimensión permite evaluar las condiciones institucionales necesarias para ejecutar las políticas de manera coherente (\citealp{oecd2024sistema,munive2022integridad}).

\textbf{Segunda dimensión: políticas de integridad y prevención.} Incluye transparencia, ética pública, gestión de conflictos de intereses, gestión de riesgos, cultura organizacional, sistemas de integridad y mecanismos de denuncia. Su función analítica es identificar cómo las instituciones buscan reducir riesgos y fortalecer conductas compatibles con el interés público (\citealp{pcm2017politica,oecd2017integrity}).

\textbf{Tercera dimensión: control, rendición de cuentas y sanción.} Comprende control interno, supervisión, rendición de cuentas, sistemas disciplinarios, investigación y sanción. Permite analizar la capacidad del sistema para detectar incumplimientos, establecer responsabilidades y evitar que los controles funcionen únicamente como procedimientos formales (\citealp{oecd2017integrity,pcm2017politica}).

El marco supone una relación secuencial y circular: las capacidades institucionales permiten implementar políticas preventivas; estas políticas reducen y gestionan riesgos; y los mecanismos de control y sanción generan retroalimentación para corregir debilidades institucionales. En conjunto, las tres dimensiones permiten estudiar la prevención de la corrupción como una función integral de la gobernanza pública.

%% ----------------------------------------------------------
\section{Metodología}
\subsection{Diseño de investigación}
El estudio se plantea como una revisión documental narrativa, con enfoque cualitativo y orientación analítica. El corpus está constituido exclusivamente por cinco documentos definidos de antemano por el investigador: tres publicaciones de la OCDE, un artículo académico que propone un Sistema Nacional de Integridad y la Política Nacional de Integridad y Lucha contra la Corrupción de la Presidencia del Consejo de Ministros del Perú (\citealp{oecd2017integrity,oecd2024justice,oecd2024sistema,munive2022integridad,pcm2017politica}). El propósito no es estimar estadísticamente la corrupción, sino integrar conceptualmente los principales planteamientos sobre gobernanza, integridad, prevención y control en el Perú.

\subsection{Fuentes y estrategia de búsqueda}
Las fuentes corresponden a documentos institucionales y académicos directamente relacionados con integridad pública, gobernanza, justicia y lucha contra la corrupción en el Perú. La estrategia de búsqueda propuesta consiste en localizar y verificar los documentos a partir de sus títulos, autores institucionales o personales y años de publicación, utilizando preferentemente los repositorios oficiales de la OCDE, la Presidencia del Consejo de Ministros y la publicación académica correspondiente.

No se proporcionaron bases de datos bibliográficas, cadenas de búsqueda ni fechas específicas de búsqueda. El investigador debe verificar estos datos antes de presentar la versión metodológica definitiva.

\subsection{Criterios de inclusión y exclusión}
Los criterios de inclusión son: documentos limitados a las cinco fuentes definidas; textos centrados en el Perú; publicaciones relacionadas directamente con gobernanza, integridad pública, prevención de la corrupción, capacidades institucionales, control o sanción; y documentos con contenido suficiente para identificar conceptos, propuestas institucionales o mecanismos de política pública.

Se excluyen fuentes adicionales, estudios comparativos ajenos al corpus definido, estadísticas externas, otros artículos académicos y documentos cuya atención principal recaiga en otros países o ámbitos no relacionados con la gestión pública peruana. Estos criterios mantienen la revisión dentro del alcance documental establecido.

\subsection{Procedimiento de análisis}
El análisis propuesto comprende cuatro etapas. Primero, se identifican los conceptos y argumentos centrales de cada fuente. Segundo, se organiza el material según tres dimensiones: coordinación y capacidades institucionales; políticas de integridad y prevención; y mecanismos de control, rendición de cuentas y sanción. Tercero, se contrastan las coincidencias y diferencias entre los documentos para establecer relaciones conceptuales entre las dimensiones. Finalmente, se elabora una síntesis interpretativa guiada por la pregunta de investigación y el marco analítico.

El procedimiento busca evitar una mera descripción individual de las fuentes y, en cambio, establecer conexiones entre arquitectura institucional, prevención y control. No se proporcionaron una matriz de extracción, un protocolo de codificación ni un procedimiento de doble revisión. El investigador debe verificar estos elementos e incorporarlos a la metodología definitiva cuando corresponda.

\subsection{Limitaciones}
La primera limitación deriva del tamaño y la composición del corpus documental, deliberadamente restringido a cinco fuentes. Por ello, la revisión no pretende representar toda la producción académica sobre corrupción e integridad en el Perú.

La segunda limitación se relaciona con el diseño narrativo, que prioriza la integración conceptual por encima de la medición cuantitativa de resultados. Por tanto, las relaciones planteadas en el marco analítico deben entenderse como categorías de interpretación y no como relaciones causales demostradas estadísticamente.

Finalmente, la falta de información sobre las bases de datos consultadas, las fechas y los términos de búsqueda y los procedimientos formales de selección impide afirmar que se aplicó una estrategia bibliométrica o una revisión sistemática. El investigador debe completar y verificar estos datos antes de presentar el artículo definitivo.

%% ----------------------------------------------------------
\section{Hallazgos}
\subsection{Responsabilidades institucionales y coordinación}
\relleno{Hallazgos pendientes de incorporar.}

\subsection{Políticas de integridad e implementación}
\relleno{Hallazgos pendientes de incorporar.}

\subsection{Mecanismos de prevención y control}
\relleno{Hallazgos pendientes de incorporar.}

\subsection{Limitaciones reportadas y brechas institucionales}
\relleno{Hallazgos pendientes de incorporar.}

%% ----------------------------------------------------------
\section{Discusión}
\subsection{Hallazgos en relación con la literatura y el marco analítico}
\relleno{Interpreta los hallazgos en relación con la literatura revisada y explica cómo responden a la pregunta de investigación.}

\subsection{Implicaciones para las políticas públicas y la gestión pública}
\relleno{Analiza las implicaciones derivadas de los hallazgos.}

\subsection{Condiciones, tensiones y desafíos pendientes}
\relleno{Analiza las condiciones, tensiones o desafíos pendientes pertinentes.}

% Ejemplo de tabla (con título numerado; no uses colores para transmitir significado):
% \begin{table}[h]
% \small\sffamily
% \caption{Título de la tabla.}
% \begin{tabular}{lll}
% \hline
% Dimensión & Hallazgo & Fuente \\
% \hline
% \relleno{texto} & \relleno{texto} & \relleno{clave} \\
% \hline
% \end{tabular}
% \end{table}

%% ----------------------------------------------------------
\section{Conclusiones y recomendaciones}
\relleno{Responde la pregunta de investigación, resume la contribución, presenta recomendaciones basadas en la evidencia, reconoce las limitaciones e identifica futuras investigaciones. No introduzcas hallazgos nuevos.}

\bibliographystyle{SageH}
\bibliography{referencias}

\end{document}
